\documentclass[12pt]{article}
\usepackage[utf8]{inputenc}
\usepackage[T1]{fontenc}
\usepackage{hyperref}
\title{Tavoyan Language Resources for K-12 Curriculum Development}
\author{Compiled by Claude Sonnet 5 (claude-sonnet-5) for Andr\'as Kornai}
\date{}
\begin{document}
\maketitle

\section{Introduction}
Tavoyan, or Dawei (ISO 639-3 tvn; Glottocode tavo1242), is a Sino-Tibetan variety closely related to Burmese. It is spoken around Dawei (Tavoy) in the Tanintharyi Region of southern Myanmar. Speakers identify as Bamar. Speaker estimates vary widely, from about 400,000 to over a million. Whether it is a separate language or a Burmese dialect is debated.

\section{Primary Educational Resources}
None found. Standard Burmese is the school language.

\section{Linguistic and Reference Materials}
\begin{itemize}
\item Overview article on Tavoyan dialects: \url{https://en.wikipedia.org/wiki/Tavoyan_dialects}
\item Academic paper (Kyoto University, Sumida) turned up in the search; its content was not verified: \url{https://www.asafas.kyoto-u.ac.jp/dl/publications/no_2002/02_sumida.pdf}
\end{itemize}

\section{Religious and Community Resources}
People-group profiles from mission-oriented organisations:
\begin{itemize}
\item \url{https://joshuaproject.net/index.php/people_groups/19250/BM}
\item \url{https://www.operation.asiaharvest.org/MYANMAR-Dawei.pdf}
\end{itemize}

\section{Digital and Audio Resources}
None found.

\section{Curriculum Development Context}
No Tavoyan-specific teaching materials were found. Tavoyan keeps features that standard Burmese has lost, such as the /-l-/ medial, and has different vocabulary, including Malay and Thai loanwords. A curriculum could use Burmese materials and add a Tavoyan component.
\end{document}